% Source: https://artofproblemsolving.com/community/c3h3146489_2023_bmmt_team_round__berkley_mini_math_tournament_fall
% Stable identity: aops-topic-3146489-q18
Every complete block $2023$ is divisible by $7$. A prefix of length $4k+r$, with $r=0,1,2,3$, is therefore congruent modulo $7$ to $0,2,20,202$, respectively. Only lengths divisible by four can give a multiple of $7$, and those prefixes end in $3$, so they are not divisible by $2,4,5,6,8,10,12$. Hence no prefix is divisible by all twelve integers, and a prefix divisible by eleven must miss only $7$ and must end in $0$. Such a prefix has length $4k+2$. Its last three digits are $320$ when $k\ge1$, so it is divisible by $8$ and $5$, hence by $2,4,10$. Its digit sum is $7k+2$, so divisibility by $9$ requires $k\equiv1\pmod9$ and then also gives divisibility by $3,6,12$. Its alternating digit sum is $k+2$, so divisibility by $11$ requires $k\equiv9\pmod{11}$. Write $k=1+9j$; then $9j\equiv8\pmod{11}$, giving $j\equiv7\pmod{11}$. The smallest nonnegative solution is $k=64$. This prefix is divisible by all integers from $1$ to $12$ except $7$, reaching the maximum eleven, and the first such length is $4\cdot64+2=\boxed{258}$.
